% This file is part of the modular source for Stratified Relational Workspace Theory.
% !TeX root = ../Stratified_Relational_Workspace_Theory.tex
\theorypart{III}{Exact Shared Cores}

\topic{6}{Global Compatible Sections}

From this point through Theorem~\ref{thm:limit-holonomy}, assume the core maps
define a local system on a connected graph unless a weaker hypothesis is
stated explicitly.

\begin{definition}[Global compatible section]\label{def:global-section}
The exact structural workspace is
\[
\boxed{
\Gamma(H)
=
\left\{
(c_i)_{i\in V}\in\bigoplus_{i\in V}C_i:
\ c_{t(e)}=H_ec_{s(e)}\ \forall e\in E
\right\}.
}
\]
An element \(c=(c_i)\in\Gamma(H)\) is a global compatible section.
\end{definition}

Equivalently,
\[
\Gamma(H)=\varprojlim_{\Path(G)}C_i.
\]
A section is not a single vector copied verbatim at every vertex. It is one
abstract element represented by a family of local vectors that agree under the
specified translation maps.

A nontrivial exact structural workspace exists when
\[
\dim\Gamma(H)>0.
\]
Its exact dimension is
\[
d_{\Workspace}=\dim\Gamma(H).
\]

Fix a root \(o\). The evaluation map is
\[
\ev_o:\Gamma(H)\to C_o,
\qquad
\ev_o((c_i))=c_o.
\]

\begin{lemma}[Root determination]\label{lem:root-determination}
If \(|G|\) is connected and every core transfer is invertible, then
\(\ev_o\) is injective. Consequently,
\[
\dim\Gamma(H)\leq\dim C_o.
\]
\end{lemma}

\begin{proof}
Choose for each vertex \(i\) a path \(\tau_i:o\to i\). Compatibility gives
\[
c_i=H_{\tau_i}c_o.
\]
Thus a section with \(c_o=0\) vanishes at every vertex. Hence \(\ev_o\) is
injective.
\end{proof}

\begin{lemma}[Tree propagation]\label{lem:tree-propagation}
If \(|G|\) is a tree and every \(H_e\) is invertible, then every
\(x\in C_o\) extends uniquely to a global compatible section.
\end{lemma}

\begin{proof}
A tree contains a unique undirected path \(\tau_i:o\to i\) for every vertex.
Define \(c_i=H_{\tau_i}x\). Because every edge belongs to these unique path
relations, the family is compatible. Uniqueness follows from
Lemma~\ref{lem:root-determination}.
\end{proof}

The obstruction to a global core therefore appears only when distinct paths
must agree. Cycles encode exactly those consistency constraints.

\topic{7}{Holonomy and Path Obstruction}

Fix a root \(o\in V\). For every loop class
\[
[\gamma]\in\pi_1(|G|,o),
\]
define the holonomy
\[
\Hol([\gamma])=H_\gamma\in GL(C_o).
\]
Functoriality gives a group representation
\[
\operatorname{hol}_o:
\pi_1(|G|,o)\longrightarrow GL(C_o).
\]
Its image is the holonomy group
\[
\Hol_o=\operatorname{im}(\operatorname{hol}_o).
\]

The obstruction associated with a loop is
\[
\Omega_\gamma=H_\gamma-I_{C_o}.
\]
The maximal path-independent root core is
\[
\boxed{
C_*
=
\bigcap_{\gamma:o\to o}\ker(H_\gamma-I).
}
\]
Equivalently,
\[
C_*=C_o^{\Hol_o}=\operatorname{Fix}(\Hol_o).
\]

\subtopic{7.1 Spanning-tree reduction}

Choose a spanning tree \(\mathcal T\subseteq|G|\). Let
\[
\tau_i:o\to i
\]
be the unique tree path. Every non-tree edge \(e:i\to j\) determines a
fundamental loop
\[
\gamma_e=\tau_j^{-1}e\tau_i
\]
with transfer
\[
H_{\gamma_e}=H_{\tau_j}^{-1}H_eH_{\tau_i}.
\]
The fundamental group of a finite graph is generated freely by these cotree
loops. Hence it is enough to impose fixed-point conditions for a finite set of
generators.

\begin{lemma}[Finite obstruction set]\label{lem:finite-obstruction}
Let \(e_1,\ldots,e_r\) be the non-tree edges of a connected finite graph.
Then
\[
C_*
=
\bigcap_{q=1}^{r}
\ker(H_{\gamma_{e_q}}-I).
\]
\end{lemma}

\begin{proof}
A vector fixed by every generator is fixed by each inverse generator and by
every finite product of generators and inverses. Those products represent all
elements of \(\pi_1(|G|,o)\). The reverse inclusion is immediate.
\end{proof}

The number of independent graph cycles is
\[
r=|E_{\mathrm{und}}|-|V|+1
\]
for a connected underlying graph, after treating parallel edges with their
multiplicity.

\topic{8}{Exact Existence Theorem}

\begin{theorem}[Limit--holonomy equivalence]\label{thm:limit-holonomy}
Let \(|G|\) be connected and let
\[
H:\Pi_1(G)\to\mathbf{FinHilb}^{\times}
\]
be a finite-dimensional local system. For any root \(o\in V\), evaluation at
the root induces a canonical vector-space isomorphism
\[
\boxed{
\ev_o:\Gamma(H)
\xrightarrow{\ \cong\ }
C_o^{\Hol_o}
=
\bigcap_{\gamma:o\to o}\ker(H_\gamma-I).
}
\]
Consequently,
\[
\Gamma(H)\neq\{0\}
\]
if and only if the holonomy representation contains a nonzero trivial
subrepresentation.
\end{theorem}

\begin{proof}
Let \(c=(c_i)\in\Gamma(H)\). For any loop \(\gamma:o\to o\), compatibility
along the loop gives
\[
H_\gamma c_o=c_o.
\]
Therefore \(\ev_o(c)=c_o\in C_o^{\Hol_o}\). By
Lemma~\ref{lem:root-determination}, \(\ev_o\) is injective.

Conversely, let \(x\in C_o^{\Hol_o}\). Choose a path
\(\tau_i:o\to i\) for each vertex and define
\[
c_i=H_{\tau_i}x.
\]
This definition is independent of the chosen path. Indeed, if
\(\sigma_i:o\to i\) is another path, then
\[
\sigma_i^{-1}\tau_i:o\to o
\]
is a loop, and invariance of \(x\) gives
\[
H_{\sigma_i}^{-1}H_{\tau_i}x=x.
\]
Multiplying by \(H_{\sigma_i}\) yields
\[
H_{\tau_i}x=H_{\sigma_i}x.
\]
For any edge \(e:i\to j\), the paths \(e\tau_i\) and \(\tau_j\) both run
from \(o\) to \(j\). Path independence therefore gives
\[
H_ec_i
=H_eH_{\tau_i}x
=H_{\tau_j}x
=c_j.
\]
Thus \((c_i)\) is a global compatible section and \(\ev_o((c_i))=x\).
Hence \(\ev_o\) is surjective and therefore an isomorphism.
\end{proof}

\begin{corollary}[Tree case]\label{cor:tree-core}
If \(|G|\) is a tree, then \(\Hol_o=\{I\}\) and
\[
\Gamma(H)\cong C_o.
\]
In particular,
\[
d_{\Workspace}=\dim C_o.
\]
\end{corollary}

\begin{corollary}[Finite matrix test]\label{cor:finite-matrix-test}
Let \(\gamma_1,\ldots,\gamma_r\) be fundamental cycle generators and suppose
\(\dim C_o=d\). Define the stacked obstruction matrix
\[
\mathcal M
=
\begin{pmatrix}
H_{\gamma_1}-I\\
\vdots\\
H_{\gamma_r}-I
\end{pmatrix}.
\]
Then
\[
\boxed{
C_*=\ker\mathcal M,
\qquad
d_{\Workspace}=d-\operatorname{rank}(\mathcal M).
}
\]
\end{corollary}

\begin{proof}
The kernel of a vertically stacked operator is the intersection of the kernels
of its blocks. Apply Lemma~\ref{lem:finite-obstruction} and
Theorem~\ref{thm:limit-holonomy}.
\end{proof}

\begin{corollary}[Averaging projector for finite holonomy]\label{cor:averaging-projector}
If \(\Hol_o\) is finite and acts unitarily, then
\[
P_*=
\frac{1}{|\Hol_o|}
\sum_{h\in\Hol_o}h
\]
is the orthogonal projector onto \(C_*\).
\end{corollary}

\begin{proof}
For every \(g\in\Hol_o\), left multiplication permutes the group elements, so
\(gP_*=P_*\). Thus \(\operatorname{ran}P_*\subseteq C_*\). If
\(x\in C_*\), then \(P_*x=x\), so \(C_*\subseteq\operatorname{ran}P_*\).
Unitarity gives \(P_*^*=P_*\), and the range characterization implies
\(P_*^2=P_*\).
\end{proof}

\subtopic{8.1 Worked examples}

\begin{example}[One-dimensional core surviving a contradictory cycle]\label{ex:reflection-cycle}
Let a three-node cycle have \(C_1=C_2=C_3=\mathbb R^2\) and
\[
H_{12}=I,
\qquad
H_{23}=I,
\qquad
H_{31}=J,
\qquad
J=
\begin{pmatrix}
1&0\\
0&-1
\end{pmatrix}.
\]
The cycle holonomy is \(J\), hence
\[
C_*=\ker(J-I)=\operatorname{span}\{(1,0)^\top\}.
\]
The first coordinate circulates consistently through all three strata. The
second returns with its sign reversed and therefore cannot define a nonzero
exact shared section. Thus
\[
d_{\Workspace}=1.
\]
\end{example}

\begin{example}[Local compatibility without a global nonzero core]\label{ex:rotation-obstruction}
Let the only nontrivial cycle holonomy on \(\mathbb R^2\) be a rotation
\[
R_\theta=
\begin{pmatrix}
\cos\theta&-\sin\theta\\
\sin\theta&\cos\theta
\end{pmatrix}
\]
with \(\theta\notin2\pi\mathbb Z\). Then
\[
\ker(R_\theta-I)=\{0\}.
\]
Every edge map is individually invertible and well-defined, yet the complete
cycle admits no nonzero exact shared core. Pairwise translatability is
therefore insufficient for global path independence.
\end{example}

\begin{example}[Acyclic collaboration]\label{ex:tree-workspace}
Suppose a root concept is translated into a proof, implementation, and report
along a tree of invertible exact transfers. No alternative route returns a
conflicting representation. The tree has trivial holonomy, so every root-core
direction extends uniquely. This does not yet establish causal efficacy; it
establishes only exact structural compatibility.
\end{example}

\subtopic{8.2 Interpretive consequence}

The theorem gives a precise meaning to ``something shared across
representations.'' The shared object is neither any one local vector nor an
extra vector outside the system. It is the equivalence class of all compatible
local realizations, canonically represented at a root by the subspace fixed by
every translation cycle.

\topic{9}{Operator Closure Theorem}

\begin{theorem}[Global operation closure]\label{thm:operator-closure}
Suppose every core transfer intertwines the representations of
\(\mathcal A\):
\[
H_e\rho_{s(e)}(a)=\rho_{t(e)}(a)H_e
\qquad
\forall e\in E,\ a\in\mathcal A.
\]
Then \(\Gamma(H)\) is an \(\mathcal A\)-module under the componentwise action
\[
\rho_\Gamma(a)(c_i)_{i\in V}
=
(\rho_i(a)c_i)_{i\in V}.
\]
Equivalently, the holonomy-fixed root core is invariant:
\[
\boxed{
\rho_o(a)C_*\subseteq C_*
\qquad
\forall a\in\mathcal A.
}
\]
\end{theorem}

\begin{proof}
Let \(c=(c_i)\in\Gamma(H)\). For every edge \(e:i\to j\),
\[
H_e\rho_i(a)c_i
=
\rho_j(a)H_ec_i
=
\rho_j(a)c_j.
\]
Therefore the family \((\rho_i(a)c_i)\) is compatible and belongs to
\(\Gamma(H)\). The representation identities on each \(C_i\) hold
componentwise, so \(\rho_\Gamma\) is an \(\mathcal A\)-representation.

Under the isomorphism of Theorem~\ref{thm:limit-holonomy}, this action is
represented at the root by \(\rho_o\). Alternatively, if \(x\in C_*\), then
Corollary~\ref{cor:holonomy-commutes} gives, for every loop \(\gamma\),
\[
H_\gamma\rho_o(a)x
=
\rho_o(a)H_\gamma x
=
\rho_o(a)x.
\]
Hence \(\rho_o(a)x\in C_*\).
\end{proof}

\begin{corollary}[Inherited algebraic structure]\label{cor:inherited-structure}
Any algebraic relation satisfied by \(\rho_i(\mathcal A)\) is inherited by the
global action. In particular, if \(a,b\in\mathcal A\), then
\[
\rho_\Gamma(ab)=\rho_\Gamma(a)\rho_\Gamma(b),
\qquad
\rho_\Gamma(1)=I_{\Gamma(H)}.
\]
\end{corollary}

The result distinguishes a shared \emph{operational} representation from a
shared static code. A core closed under the common algebra supports the same
abstract manipulations across strata without leaving the globally compatible
subspace.

\topic{10}{Workspace Operator Algebra}

\begin{definition}[Natural workspace operators]\label{def:workspace-operator-algebra}
Define
\[
\boxed{
\mathcal O_{\Workspace}
=
\left\{
(A_i)_{i\in V}:
A_{t(e)}H_e=H_eA_{s(e)}\ \forall e
\right\}.
}
\]
Categorically,
\[
\mathcal O_{\Workspace}=\Nat(H,H),
\]
the algebra of natural endomorphisms of the local system.
\end{definition}

Pointwise addition, scalar multiplication, and composition make
\(\mathcal O_{\Workspace}\) a unital associative algebra. Each
\(A=(A_i)\) acts on global sections by
\[
\pi_{\Workspace}(A)(c_i)=(A_ic_i).
\]
Naturality guarantees that \((A_ic_i)\) is again compatible.

\begin{theorem}[Commutant characterization]\label{thm:commutant-characterization}
Let \(|G|\) be connected and let \(H\) be an invertible local system. Root
evaluation induces an algebra isomorphism
\[
\boxed{
\Phi_o:\mathcal O_{\Workspace}
\xrightarrow{\ \cong\ }
\operatorname{Comm}(\Hol_o),
\qquad
\Phi_o((A_i))=A_o,
}
\]
where
\[
\operatorname{Comm}(\Hol_o)
=
\{B\in\End(C_o):Bh=hB\ \forall h\in\Hol_o\}.
\]
\end{theorem}

\begin{proof}
If \((A_i)\in\mathcal O_{\Workspace}\), naturality along every loop gives
\[
A_oH_\gamma=H_\gamma A_o.
\]
Thus \(A_o\in\operatorname{Comm}(\Hol_o)\), and \(\Phi_o\) is a well-defined
algebra homomorphism.

Choose paths \(\tau_i:o\to i\). Naturality along \(\tau_i\) gives
\[
A_iH_{\tau_i}=H_{\tau_i}A_o,
\]
so
\[
A_i=H_{\tau_i}A_oH_{\tau_i}^{-1}.
\]
Therefore \(A_o\) determines every \(A_i\), and \(\Phi_o\) is injective.

Conversely, take \(B\in\operatorname{Comm}(\Hol_o)\) and define
\[
A_i=H_{\tau_i}BH_{\tau_i}^{-1}.
\]
This is independent of the chosen root path: changing \(\tau_i\) multiplies
it by a loop holonomy that commutes with \(B\). For an edge \(e:i\to j\), the
operator
\[
H_{\tau_j}^{-1}H_eH_{\tau_i}
\]
is a root holonomy. Commutation with \(B\) rearranges to
\[
A_jH_e=H_eA_i.
\]
Hence \((A_i)\) is natural and \(\Phi_o((A_i))=B\). Thus \(\Phi_o\) is
surjective.
\end{proof}

The structural workspace is therefore more accurately represented by the
triple
\[
\boxed{
\Workspace_{\mathrm{static}}
=
\bigl(
\Gamma(H),
\mathcal O_{\Workspace},
\pi_{\Workspace}
\bigr),
}
\]
not merely by the vector space \(\Gamma(H)\).

\subtopic{10.1 Effective operator algebra}

Some natural operators may vanish on every global section while acting
nontrivially on locally available but globally inconsistent directions. Define
the annihilator
\[
\operatorname{Ann}(\Gamma(H))
=
\{A\in\mathcal O_{\Workspace}:
\pi_{\Workspace}(A)=0\}.
\]
The effective workspace algebra is
\[
\boxed{
\mathcal O_{\Workspace}^{\mathrm{eff}}
=
\mathcal O_{\Workspace}/\operatorname{Ann}(\Gamma(H)).
}
\]
Its action on \(\Gamma(H)\) is faithful by construction.

\begin{example}[Operators in the reflection-cycle workspace]\label{ex:reflection-operators}
For Example~\ref{ex:reflection-cycle}, the holonomy is
\[
J=\operatorname{diag}(1,-1).
\]
Its commutant consists of diagonal matrices:
\[
\operatorname{Comm}(J)
=
\left\{
\begin{pmatrix}a&0\\0&b\end{pmatrix}:a,b\in\mathbb R
\right\}.
\]
The exact workspace core is the first coordinate. Consequently, only the
parameter \(a\) acts effectively on \(\Gamma(H)\); the \(b\)-direction belongs
to the annihilator of the global core. Hence
\[
\mathcal O_{\Workspace}^{\mathrm{eff}}\cong\mathbb R.
\]
\end{example}

\subtopic{10.2 Scope of the exact theory}

Parts~I--III now establish three exact facts:
\[
\begin{aligned}
&\text{global shared content is the holonomy-fixed core},\\
&\text{universally transported operations preserve that core},\\
&\text{natural workspace operations are the commutant of holonomy}.
\end{aligned}
\]
These statements do not yet prove that the shared core is robust, causally
used, irreducibly distributed, or temporally persistent. Those are separate
questions developed in the approximate and dynamic theory that follows.
